\unit{Fine structure, hyperfine structure and fields}
\unitref{Weeks 8--9; Assignment 3 Q4; 2025 Q1}
\subsection{Why one electronic configuration has many levels}
The central-field approximation produces
orbital energies but does not yet include
all small interactions. Electron spin--orbit
coupling links the orbital and spin
magnetic moments; the atomic nucleus has
its own spin and magnetic moment; external
electric and magnetic fields shift the levels.
These effects create fine, hyperfine,
Zeeman and Stark structures with
different typical energy scales.

\subsection{Fine structure from \(\mathbf L\cdot\mathbf S\)}
Suppose \(H_{\rm SO}=A_{\rm SO}\,
\mathbf L\cdot\mathbf S/\hbar^2\)
within an LS term. Since
\[
\mathbf J^2=(\mathbf L+\mathbf S)^2
=\mathbf L^2+\mathbf S^2+
2\mathbf L\cdot\mathbf S,
\]
the energy correction is
\[
\boxed{E_J^{\rm SO}=\frac{A_{\rm SO}}2
[J(J+1)-L(L+1)-S(S+1)].}
\]
This is more than a formula to remember:
it is simply the scalar product identity
for two angular momenta.

\begin{worked}{Split a \({}^3D\) term}
The term has \(S=1,L=2\), so
\(J=1,2,3\). With \(L(L+1)=6\) and
\(S(S+1)=2\),
\[
E_1=-3A_{\rm SO},\quad
E_2=-A_{\rm SO},\quad
E_3=2A_{\rm SO}.
\]
The degeneracies \(3+5+7=15\)
equal \((2L+1)(2S+1)=5\times3\).
The level intervals \(2A_{\rm SO}\)
and \(3A_{\rm SO}\) are not equally spaced.
\end{worked}

\subsection{Hyperfine structure couples electrons to the nucleus}
The nuclear spin \(\mathbf I\) combines
with electronic \(\mathbf J\) to form
\(\mathbf F=\mathbf I+\mathbf J\).
For magnetic-dipole hyperfine coupling,
\[
H_{\rm hfs}=\frac{A}{\hbar^2}\mathbf I\cdot\mathbf J,
\]
giving
\[
\boxed{E_F=\frac A2[
F(F+1)-I(I+1)-J(J+1)]},\quad
F=|I-J|,\ldots,I+J.
\]
The numerical sign of \(A\) matters:
the labels alone do not establish
which \(F\) has higher energy.
Electric quadrupole and higher-rank
nuclear hyperfine interactions are
possible only if the nuclear and
electronic tensor ranks allow them.

\begin{worked}{The 2025 exam hyperfine example}
Let \(I=5/2,J=1/2\).
Then \(F=2,3\).
Since \(I(I+1)=35/4\),
\(J(J+1)=3/4\),
\[
E_2=\frac A2(6-38/4)=
\boxed{-7A/4},
\quad
E_3=\frac A2(12-38/4)=
\boxed{5A/4}.
\]
Their separation is \(3|A|\).
The magnetic substate count
\((2F_1+1)+(2F_2+1)=5+7=12\)
agrees with \((2I+1)(2J+1)=6\times2\).
\end{worked}

\subsection{The weak magnetic-field regime}
A magnetic field interacts with atomic
magnetic moments. When hyperfine coupling
dominates the Zeeman energy,
\(F,m_F\) remain useful approximate
quantum numbers:
\[
\boxed{\Delta E_{Fm_F}\simeq
\mu_Bg_Fm_FB.}
\]
Neglecting the much smaller nuclear
magnetic contribution and for \(F>0\),
\[
g_F\simeq g_J
\frac{F(F+1)+J(J+1)-I(I+1)}
{2F(F+1)}.
\]
The electronic Land\'e factor for
\(g_L=1,g_S\simeq2\) is
\[
\boxed{g_J\simeq1+
\frac{J(J+1)+S(S+1)-L(L+1)}
{2J(J+1)}}.
\]
These factors arise from \emph{projecting}
electronic orbital/spin moments first
onto \(J\), then onto \(F\).

\subsection{The strong-field (Paschen--Back) regime}
When Zeeman energy exceeds hyperfine
coupling, nuclear and electronic
angular momenta approximately decouple.
The natural basis changes from
\(\ket{Fm_F}\) to \(\ket{Jm_J,Im_I}\).
To first order,
\[
E\simeq g_J\mu_Bm_JB+Am_Im_J
\]
if the nuclear Zeeman contribution is
neglected. The total state count stays
constant, but the interpretation
of the levels changes.
A fixed-\(F\) weak-field formula
must \emph{not} be extrapolated
to arbitrarily large \(B\).

\subsection{Electric fields: the Stark effect}
A static electric field \(\mathcal E\hat z\)
couples through \(-\mathbf d\cdot\boldsymbol{\mathcal E}\).
For a nondegenerate parity eigenstate
with no permanent electric dipole,
the first-order shift vanishes.
Second-order perturbation theory gives
\[
\boxed{\Delta E^{(2)}=-\frac12
\alpha(0)\mathcal E^2,\quad
\alpha(0)=2\sum_{n\ne i}
\frac{|\langle n|d_z|i\rangle|^2}{E_n-E_i}}.
\]
The polarisability measures how readily
the atom's electron cloud responds to
an applied field. For a hydrogenic \(1s\)
state a closure approximation using
\(E_n-E_1\approx E_2-E_1\) gives
\(\alpha_{\rm approx}=16/(3Z^4)\)
in atomic polarizability-volume units,
compared with the exact hydrogen value
\(9/2\) at \(Z=1\).

\begin{checkpoint}
Identify the relevant coupling hierarchy.
Fine structure: \(L+S\to J\).
Hyperfine: \(I+J\to F\).
Weak magnetic field: \(F,m_F\).
Strong magnetic field: \(m_I,m_J\).
Electric field: check parity before
applying Stark perturbation theory.
\end{checkpoint}
